% TOP_PROBLEM: {"id": 3700023, "title": "Equilibrium for infinitely many positive masses", "category_id": 8, "category": {"id": 8, "name": "analysis", "display_name": "Analysis", "description": "Limits, continuity, calculus, and function theory.", "slug": "analysis", "order_index": 8, "created_at": "2026-07-31T15:26:25.671Z"}, "status": "partially_solved", "problem_number": "AMR-036-0023", "set_id": 13, "statusQualification": "Restores the primary source's essential infinite-configuration hypothesis and dimension n>=2; uses an equivalent tail summability convention permitting an origin mass."}
% FIRST_POSTED: 2026-10-01
% ENTRY_KIND: statement_only
% UnsolvedMath ID 3700023; source catalogue code AMR-036-0023
% !TeX program = lualatex
% Definition and sources only; this file is not an LLM solution attempt.
\documentclass[11pt,a4paper]{article}
\usepackage[margin=25mm]{geometry}
\usepackage{fontspec}
\setmainfont{lmroman10-regular.otf}[BoldFont=lmroman10-bold.otf,
 ItalicFont=lmroman10-italic.otf,BoldItalicFont=lmroman10-bolditalic.otf]
\usepackage{amsmath,amssymb,mathtools}
\usepackage{unicode-math}
\setmathfont{latinmodern-math.otf}
\AtBeginDocument{\let\setminus\smallsetminus}
\usepackage{xurl,hyperref}
\hypersetup{colorlinks=true,urlcolor=blue}
\setcounter{secnumdepth}{0}
\setlength{\parindent}{0pt}
\setlength{\parskip}{5pt}
\setlength{\emergencystretch}{3em}
\begin{document}
\section*{Equilibrium for infinitely many positive masses}\label{3700023}
{\small UnsolvedMath ID 3700023 · AMR-036-0023 · Analysis · AMR Open Problem Lists}
\subsection{Definitions and mathematical statement}
Fix a dimension $n\geq2$. Let $(x_k)_{k\geq1}$ be an
infinite sequence of distinct points of $\mathbb R^n$
with no finite accumulation point, and let $a_k>0$ be
their masses. Assume the convergence condition
\[
             \sum_{k=1}^{\infty}
                    \frac{a_k}{(1+|x_k|)^{n-1}}<\infty.
\]
The use of $1+|x_k|$ allows a mass at the origin; outside
a fixed ball this is equivalent to the denominator in
the source formulation.

For $x$ outside the set of masses define the force
\[
       F(x)=\sum_{k=1}^{\infty}
                 \frac{a_k(x_k-x)}{|x-x_k|^n}.
\]
The displayed summability condition makes this vector series
absolutely and locally uniformly convergent away from the
mass points: its tail is bounded by a constant multiple
of $\sum_k a_k/|x_k|^{n-1}$ on each fixed compact set.

Must there exist a point
$x\in\mathbb R^n\setminus\{x_k:k\geq1\}$ with $F(x)=0$?
Such a point is an equilibrium of the logarithmic force
when $n=2$ and of the Newtonian force when $n\geq3$.
The question concerns an infinite discrete configuration
and allows masses tending to zero, arbitrary geometry of
their positions, and infinite total mass whenever the
weighted sum converges. No uniform positive lower bound
on the $a_k$ and no stronger convergence assumption for
an unrenormalized potential is imposed.
\subsection{Short English statement}
Must the convergent force field of infinitely many positive, locally discrete masses have an equilibrium point?
\subsection{Sources}
\begin{itemize}
\item[S1] ULAM.AI and the UnsolvedMath Contributors, supplied problems.json, record 3700023 (AMR-036-0023), AMR Open Problem Lists. Catalogue identity and source transcription; CC BY 4.0.
\url{https://huggingface.co/datasets/ulamai/UnsolvedMath}.
\item[S2] Eremenko - My favorite unsolved problems, Equilibrium points of logarithmic and Newtonian potentials. Original source indexed by AMR-036-0023.
\url{https://www.math.purdue.edu/~eremenko/uns1.html}.
\item[S3] A. Eremenko, Equilibrium points of logarithmic and Newtonian potentials. Author-maintained primary problem note.
\url{https://www.math.purdue.edu/~eremenko/dvi/equil.pdf}.
\end{itemize}
\end{document}
